% =====================================================================
%  Chapter 13 — The Verdict, Part I
% =====================================================================

\part{The Verdict}
\label{part:verdict}
\markboth{Part V}{The Verdict}

\chapter{The Finding}
\label{ch:verdict}

This chapter states the conclusion of the audit. It is stated in one sentence
first and then defended, because the sentence is the deliverable and the
defence is the part that should be checked.

\section{The finding}\begin{seamnote}{\textbf{The finding}}
On the evidence available as of 21 November 2026, 13:55 UTC, \textbf{the
formal status of the claim is settled and the mathematical status is not}. The
development compiles; the Lean and Mathlib communities have jointly certified,
with the endorsement of the Lean FRO and the Mathlib Initiative, that the
theorem proved is exactly Mathlib's \texttt{RiemannHypothesis}, and that the
exported proof is accepted by Lean's official kernel and by several
independently implemented checkers, two of which carry formal soundness proofs
relative to set theory. \textbf{The claim is formally certified.}

What no one has yet done is read the mathematics. No mathematician has
confirmed that the load-bearing intermediate statements are the classical
theorems they are named for; the author's own account of the method records an
a-priori ceiling at two thirds and an obstruction whose stated remedy requires
a conjecture that remains open; and the remark does not explain how either was
crossed. \textbf{The claim is formally established and mathematically
unaudited}\idx{verdict}.
\end{seamnote}

Four words in that sentence carry the weight, and it is worth being exact
about what each one rules out and what it permits. The distinction between
the second and the third is the one that the coverage of 20 and 21 November
2026 ran together, and it is the one this chapter exists to keep apart.

\begin{description}[style=nextline,leftmargin=0pt,labelwidth=0pt,labelsep=0pt,font=\sffamily\bfseries\small,
  itemsep=6pt]
\item[\textbf{Formally certified.}] This is a statement of fact about the
21 November joint statement, and it is the half of the position that the
book's earlier draft of this verdict got wrong by omission. The development
compiles; the theorem proved has been matched against Mathlib's
\texttt{RiemannHypothesis} by two independent tools; the exported proof is
accepted by Lean's kernel and by independently implemented checkers, two of
which have formal soundness proofs. Whatever else is open, \emph{that} is
closed, and it is closed by the community that owns the tooling, not by this
book.

\item[\textbf{Mathematically unaudited.}] This is the half that remains open,
and it is where the review has not happened. No mathematician has read the
intermediate statements to confirm that each is the classical theorem it is
named for. The joint statement does not claim otherwise; it says it looks
forward to ``all the activity in the mathematical community that will help
explain the new ideas in this proof'', and it does not claim that any human
has read the argument. The specific mathematical content of the claim has
not been examined by any independent party, so no probability is attached to
it by this book.

\item[\textbf{Not refuted.}] Not ``probably true'', not ``likely to be
disproved''. What the evidence does support is a statement about \emph{where
the claim stands}: it is not yet a theorem in the sense mathematicians use the
word, because the process that turns a formalised development into an accepted
theorem has one step still outstanding, and that step is a human one.

\item[\textbf{Serious.}] This is the part most likely to be under-weighted by
sceptics, and it is why the book does not use the word ``failed''. Zero
axioms, zero \texttt{sorry}, a correct reduction, a bottom theorem that is
genuinely proved, a certified formal status, and a twenty-component bundle of
separately proved theorems named after the actual results one needs: that is
not the profile of a fabrication.

\end{description}

\section{Four findings that support the finding}\begin{enumerate}
\item \textbf{The artifact is clean and the reduction is sound.} The
development contains no axioms and no admitted gaps, and its eight-link chain
from Mathlib's \texttt{RiemannHypothesis} down to a proved holomorphic
Lefschetz formula has been traced link by link. The reduction from
non-vanishing of Dirichlet $L$-functions to RH is the standard argument,
correctly executed. \hfill(Chapters~2, 7, 9)

\item \textbf{The load-bearing inputs have not been audited.} The
automorphy step is carried by the bundle \texttt{Af.Hypotheses\_2}, whose
field \texttt{newform} supplies the automorphic representation the downstream
theorem consumes. The bundle's other fields are re-derived in the source from
named published results --- which is the right shape --- but the fidelity of
those re-derivations to the published theorems is precisely the question, and
no one outside the author has examined it. \hfill(Chapter~8)

\item \textbf{The author's own account records two unresolved obstacles.} The
timeline's item 14 states that the method ``would stop at 2/3 a priori'' and
lists attempts --- a main term, a transfer of constants, killing by twisting,
``division'' --- without stating that any of them succeeded. Item 17 states
that the $H^1_o$ classes that twisting fails to kill require ``a significant
weakening of Vandiver'', which is an open problem, and its footnote is
hedged with ``maybe?''. GRH is recorded for the following day, with no
account of the transition. \hfill(Chapters~1, 4, 5)

\item \textbf{The formal review has happened and succeeded; the mathematical
review has not begun.} This is the finding most in need of restating on 21
November 2026. The development compiles, independently and reproducibly; the
axiom question is settled in the affirmative; the statement has been matched
against Mathlib's by two independent tools; and the exported proof has been
accepted by Lean's kernel and by several independently implemented checkers,
with the community certifying the result jointly and publicly. The remaining
layer --- a human statement review --- has not begun in public, and the
community's own statement says so: it looks forward to the mathematical
community explaining the new ideas, and it does not claim that any human has
read the argument. The artefact's own author explains why the prose is
missing: he has ``not mastered a massive neighbourhood around its
mathematics''. \hfill(Chapters~6, 10, 12)
\end{enumerate}

\section{What this finding is not}\begin{mythbox}{Four misreadings, and why they are wrong}
\begin{itemize}
\item \textbf{``So the proof is fake.''} No. Nothing found in this audit
refutes anything, and the artifact's quality is genuinely high. The formal
status is certified by the community that owns the tooling. The correct
statement is that its \emph{mathematical} correctness has not been
established.

\item \textbf{``So it is probably right.''} No. This book attaches no
probability to the claim. The evidence is consistent with the claim being
correct, and equally consistent with a formalisation whose deep inputs are
weaker than advertised. Distinguishing these requires the review that has not
happened.

\item \textbf{``So the certification settles it.''} No, and this is the
misreading that the book was written to prevent, in the opposite direction.
The certification settles that the development is complete, that the
statement is the right one, and that the derivation is valid. It does not
settle that the intermediate statements are true. A checker verifies that a
conclusion follows from its premises; it cannot verify that the premises
are true. The joint statement does not claim otherwise, and the community's
own notes list what a passing run still trusts, among them ``Mathlib's
definitions'' and the sandboxes.

\item \textbf{``So the 2/3 ceiling proves it is wrong.''} No. The ceiling is
what the method produced on 15 November. If the Lean development contains
something that crosses it --- which is the entire claim --- the remark's
silence is explicable as compression. The remark is an admitted ``slop''
summary by a self-declared non-expert in the relevant area, and it explicitly
tells the reader that the code is the authoritative version.

\item \textbf{``So the 2/3 ceiling proves it is right.''} No, and this is the
error the coverage made. A method that a competent specialist computed would
stop at two thirds is a method with a demonstrated ceiling, not a method with
a demonstrated crossing. The remark's silence is as consistent with a
successful crossing as with a failure, and it is equally consistent with the
third possibility --- that the formal statement proved is weaker than the
headline theorem.
\end{itemize}
\end{mythbox}

\section{The three possibilities, stated precisely}Chapter~4 named three possibilities for the 15-to-16 November gap. It is worth
restating them here in their sharpest form, because they are what a reviewer
must discriminate between, and they correspond to three different outcomes.

\begin{description}[style=nextline,leftmargin=0pt,labelwidth=0pt,labelsep=0pt,font=\sffamily\bfseries\small,
  itemsep=6pt]
\item[\textbf{A. The artifact contains the missing ingredient.}] The theta-lift
chain, or the epsilon estimate, or the hyperplane-section package, or some
part of the 175 thousand files, contains a step that the remark's compressed
timeline does not describe, and that step genuinely raises the bound from two
thirds to all of the zeros. \emph{Outcome: the claim is correct, and the
remark is simply a bad write-up.} This is a live possibility and it would be
an extraordinary result.

\item[\textbf{B. The ceiling was a cost accounting the author later
improved.}] The computation of item 14 was superseded; the remark records a
superseded state of the argument. \emph{Outcome: the claim is correct, and
the remark is out of date.} Also live, and also unremarkable.

\item[\textbf{C. The formal theorem is weaker than the headline.}] The
development proves a statement that implies, or is close to, the two-thirds
proposition of August 2026, and the top-level theorem is obtained by a step
that depends on an input whose formal content differs from its informal
description. \emph{Outcome: the claim is wrong, and the artifact is a
substantial and honest piece of work that does not do quite what its headline
says.}
\end{description}

\begin{seamnote}{How this book rates the three}
This book does not rate them. That is the point of the finding: the evidence
available does not discriminate, and a book that pretended otherwise would be
guessing.

What it can say is this. Possibilities A and B are consistent with everything
known, including the method's demonstrated quality, its August result, and the
author's candour. Possibility C is consistent with everything known too, and
is the standard failure mode of formalised mathematics, and is the only one of
the three that a compiler cannot detect.

If a reader wants a bet, the honest form of the bet is: the artifact is
genuinely substantial, the author is being unusually honest about the state of
it, and the most likely outcome of a full review is a significant
\emph{reformulation} of the claim rather than either a refutation or a
triumphal confirmation. But the bet is not the finding, and the finding does
not depend on it.
\end{seamnote}

\section{What would change the finding}Three things would, in decreasing order of speed.

\begin{enumerate}
\item \textbf{A completed statement review} of the load-bearing inputs,
conducted by somebody with the relevant expertise, published. This is the
decisive step, and Chapter~14 says what it consists of.

\item \textbf{A published paper} describing the mathematics of the argument,
with the 15-to-16 November transition explained. This would not by itself
settle the mathematics, but it would make the review possible, and it would
close the gap between the two documents that this audit could not close.

\item \textbf{An independent kernel replay} via \texttt{lean4export} and an
external kernel such as NanoDa, or a check via \texttt{lean4checker}. This
would strengthen the mechanical layer. It is already effectively done; doing
it again adds assurance, not information.
\end{enumerate}

Nothing short of the first would change the finding, and the first is a
months-long project for a specialist. This is the ordinary state of affairs
for a large formalisation, and it is not a sign of anything.

One item in that list has changed since the first draft of this chapter, and
it is the one a reader of 21 November will notice: the independent kernel
replay is no longer a thing to be done. It was done. The \texttt{nanoda},
\texttt{con-leche}, \texttt{con-ron} and \texttt{lean4lean} runs are public,
with logs, and they accepted the proof. The list of what would change the
finding has one fewer entry than it had yesterday, and the one that remains
is the one that requires a person rather than a machine.

\begin{keytakeaways}

\item The finding, as of 21 November 2026 at 13:55 UTC: \textbf{formally
certified, mathematically unaudited}.
\item The formal half is closed by the community, not by this book. The
development compiles, the statement is Mathlib's, the exported proof is
accepted by Lean's kernel and by independently implemented checkers, and the
Lean FRO and the Mathlib Initiative certify it jointly.
\item The artifact is clean, its reduction is sound, and its architecture is
appropriate. This is a serious submission, not a fabrication.
\item The unexamined part is the fidelity of the load-bearing inputs to the
published theorems they encode. They are proved inside the library, so there
is no logical gap; the open question is whether they are true.
\item The remark's own timeline records two obstacles --- a two-thirds ceiling
and a Vandiver-dependent obstruction --- and does not say how either was
crossed.
\item Three outcomes remain live: the artifact contains the missing
ingredient; the remark is out of date; or the formal theorem is weaker than
the headline. The available evidence does not discriminate between them.
\end{keytakeaways}
